% Test file to check if the technical report fits in a single CVPR column
\documentclass[10pt,twocolumn,letterpaper]{article}
\usepackage{cvpr}
\usepackage{graphicx}
\usepackage[pagebackref,breaklinks,colorlinks,allcolors=blue]{hyperref}

\def\paperID{*****}
\def\confName{CVPR}
\def\confYear{2026}

\title{Column Test}

\begin{document}

% ---- Left column: our report ----
\subsection*{GatedMemorySAM}

\textit{Jemo Maeng, Sangmin Park, Seongju Lee, Kyoobin Lee}\\
{\small \texttt{\{maengjemo, leowiu24, lsj2121\}@gm.gist.ac.kr}}\\
{\small \texttt{kyoobin@gist.ac.kr}}\\
\textit{Gwangju Institute of Science and Technology}

\vspace{0.5em}

\noindent\textbf{Method.}
We introduce GatedMemorySAM, which builds upon MemorySAM~\cite{liao2025memorysam} that repurposes SAM2's~\cite{ravi2024sam2} temporal memory attention for cross-modal fusion by treating each input modality (RGB, LiDAR, Thermal) as a separate ``frame'' in SAM2's video pipeline.
We extend this framework with two key modifications: (1) Soft MoE LoRA adaptation, and (2) quality-aware modality scoring.

For backbone adaptation, we replace standard LoRA~\cite{hu2022lora} with Soft MoE routing~\cite{puigcerver2024softmoe} over multiple LoRA experts, inserted into every attention block (Q and V projections, 48 layers total) with rank$=$4 and 3 experts.
Each spatial token is softly routed to all experts via a learned gating network, enabling per-token specialization across modalities.

For modality weighting, we introduce a CrossModalFusionHead that compares all modality features via global average pooling and a shared comparison MLP to produce per-modality softmax weights.
These weights drive two score-based fusion mechanisms: (i) memory modulation, which max-normalizes the scores and scales each modality's backbone features before they enter SAM2's memory bank, so that the best modality retains full strength while weaker ones are suppressed; and (ii) weighted mask fusion, which averages the per-modality decoder outputs using the same softmax weights.
The overall pipeline is illustrated in Figure~\ref{fig:method}.

\begin{figure}[h]
\centering
\includegraphics[width=\linewidth]{fig/challenge_fig.png}
\caption{Overview of GatedMemorySAM. Each modality is encoded by SAM2 with Soft-MoE LoRA, scored by the CrossModalFusionHead, and fused via memory modulation and weighted mask fusion.}
\label{fig:method}
\end{figure}

\noindent\textbf{Training Details.}
We train with the AdamW optimizer (lr=6e-4, weight decay=0.01) using a warmup polynomial LR schedule (10-epoch warmup, power=0.9) for 200 epochs.
The loss function is Online Hard Example Mining Cross-Entropy (OHEM-CE).
Input images are resized to 1024$\times$1024.
We use DDP training with an effective batch size of 16 across 8 GPUs with gradient accumulation.

\noindent\textbf{Night Augmentation.}
Since the test set consists entirely of nighttime images while training data is daytime, we apply an extensive nighttime simulation pipeline:
brightness darkening (range [0.03, 0.45] with 60\% dark-biased sampling), contrast reduction, gamma correction ([0.4, 0.8]), Gaussian and Poisson shot noise, cross-modal replacement (CRM, p=0.20), and PhysAug~\cite{xu2025physaug}-inspired spatial perturbations (random convolution filters and planar Fourier wave patterns, p=0.40).

\noindent\textbf{Datasets and Pretrained Weights.}
We train exclusively on the MULTIAQUA dataset provided by the challenge organizers.
The SAM2 Hiera-B+ backbone is initialized from publicly available pretrained weights.
No additional external datasets or annotations are used.

\noindent\textbf{Hardware and FPS.}
Training was conducted on 8$\times$ NVIDIA RTX 3090 GPUs (24GB each).
Inference runs on a single NVIDIA Titan RTX GPU at approximately 2.1 FPS (1024$\times$1024 input, 3 modalities processed sequentially).

\noindent\textbf{Maritime Domain Adaptations.}
The primary challenge is the extreme day-night domain gap (daytime training vs.\ nighttime-only test).
Our nighttime simulation augmentation was the single most impactful component, nearly doubling the test mIoU.

\noindent\textbf{Submission.}
Our best submission (Submission \#16710, leaderboard name: \textit{kblkblkbl}) achieved a validation mIoU of 93.29\% and a test mIoU of 70.91\%, yielding an M-score of 82.10, where M $= 0.75 \times$ mIoU$_{\mathrm{val}}$ $+ 0.25 \times$ mIoU$_{\mathrm{test}}$.

{\small
\bibliographystyle{ieeetr}
\bibliography{references}
}

\end{document}
